\chapter{Laboratorio di Algoritmi -\/ FAQs}
\hypertarget{md_FAQ}{}\label{md_FAQ}\index{Laboratorio di Algoritmi -\/ FAQs@{Laboratorio di Algoritmi -\/ FAQs}}
\label{md_FAQ_autotoc_md0}%
\Hypertarget{md_FAQ_autotoc_md0}%
 {\bfseries{Q\+: Come determino quale turno di laboratorio posso frequentare?}}

{\bfseries{A}}\+: I turni di laboratorio sono suddivisi secondo due criteri\+:


\begin{DoxyEnumerate}
\item secondo il corso di assegnazione\+:
\begin{DoxyItemize}
\item i turni 1 e 2 sono destinati agli studenti che frequentano il corso A
\item i turni 3 e 4 sono destinati agli studenti che frequentano il corso B
\end{DoxyItemize}
\item secondo la parità della cifra meno significativa della matricola dello studente\+:
\begin{DoxyItemize}
\item se la cifra è pari frequentare turni pari
\item se la cifra è dispari frequentare turni dispari
\end{DoxyItemize}
\end{DoxyEnumerate}

Riassumendo\+:
\begin{DoxyItemize}
\item studenti del corso A con la cifra meno significativa della matricola dispari saranno assegnati al turno 1
\item studenti del corso A con la cifra meno significativa della matricola pari saranno assegnati al turno 2
\item studenti del corso B con la cifra meno significativa della matricola dispari saranno assegnati al turno 3
\item studenti del corso B con la cifra meno significativa della matricola pari saranno assegnati al turno 4
\end{DoxyItemize}

{\bfseries{Q\+: Sono uno/a studente/ssa assegnato/a al turno X posso spostarmi al turno Y?}}

{\bfseries{A}}\+: Si veda quanto riportato nel README.

{\bfseries{Q\+: Sono uno/a studente/ssa che ha già preparato il laboratorio dell\textquotesingle{}anno accademico 2023/2024. Posso sostenere la prova di laboratorio discutendo gli esercizi dello scorso a.\+a.?}}

{\bfseries{A}}\+: Si veda quanto riportato nel README.

{\bfseries{Q\+: Sono uno/a studente/ssa del turno X vorrei lavorare in gruppo con uno/a studente/ssa del turno Y. Posso farlo?}}

{\bfseries{A}}\+: No.

{\bfseries{Q\+: Sono uno/a studente/ssa lavoratore/lavoratrice, è prevista qualche regola particolare per chi non frequenta?}}

{\bfseries{A}}\+: No.

{\bfseries{Q\+: Ho l\textquotesingle{}impressione che si dia molta (troppa) importanza alla parte "{}estetica"{} del codice (commenti, indentazioni, etc..). Spesso passo più tempo ad abbellire il codice che a scrivere gli algoritmi stessi.}}

{\bfseries{A}}\+: Se hai questa impressione allora vuol dire che stai svolgendo bene il laboratorio. La scrittura di codice modulare, di buona qualità, ben commentato e ben testato è un lavoro che richiede tempo e dedizione. È però un\textquotesingle{}attività fondamentale per il successo dei progetti a cui lavorerai. Quando oltre a percepirne la fatica ne percepirai anche il valore, avrai interiorizzato una buona parte di quanto si voleva trasmettere quando si è formulato il mandato degli esercizi.

{\bfseries{Q\+: E\textquotesingle{} possibile sostenere la prova di esame di laboratorio da remoto?}}

{\bfseries{A}}\+: No. 